\clearpage
\fancyhead[L]{\small COMET.}
\hypertarget{native-442-part-lower}{}
\begin{center}
{\large COMET.}

\noindent\rule{\linewidth}{0.4pt}

Göttingische gelehrte Anzeigen. No. 78. pp. 769..771. 17 May 1821.

\noindent\rule{\linewidth}{0.4pt}
\end{center}

Until now we have put off reporting all the observations of this year’s comet made at our observatory. We hoped that much later observations from elsewhere might further refine the results based on ours. Since that hope has gone unfulfilled, we now give all the observations made here.

\Needspace{14\baselineskip}
\SourceLeaf{443}
\begin{center}
\begin{tabular}{lr|r|r|r}
\multicolumn{2}{c|}{1821} & Mean time & Right ascension & Northern decl. \\
\hline
Jan. & 30 & $7^{\mathrm{u}}\;34^{\mathrm{m}}\;32^{\mathrm{s}}$ & $359^\circ\;27'\;7''$ & $16^\circ\;4'\;36''$ \\
Feb. & 3 & $7\quad3\quad56$ & $359\quad3\quad54$ & $15\quad46\quad3$ \\
& 7 & $6\quad42\quad38$ & $358\quad45\quad5$ & $15\quad29\quad49$ \\
& 9 & $6\quad42\quad30$ & $358\quad36\quad24$ & $15\quad21\quad22$ \\
& 10 & $6\quad52\quad27$ & $358\quad32\quad19$ & $16\quad17\quad50$ \\
& 11 & $7\quad12\quad0$ & $358\quad28\quad27$ & $15\quad14\quad26$ \\
Mar. & 1 & $7\quad18\quad7$ & $357\quad13\quad34$ & $14\quad8\quad37$ \\
& 5 & $7\quad5\quad28$ & $356\quad54\quad11$ & $13\quad43\quad5$ \\
\end{tabular}
\end{center}

\Needspace{14\baselineskip}
We have often shared examples of Mr. von \textsc{Staudt}’s outstanding skill in astronomical calculation. He bases the following \emph{parabolic elements} on these observations.

\begin{center}
\small
\begin{tabular}{p{.73\linewidth}r}
Time of passage through perihelion 1821 March \mbox{.\,.\,.\,.\,.\,.\,.\,.\,.} & $21.61890$ \\
Logarithm of the least distance from the Sun Meridian of Göttingen & $8.9645990$ \\
Argument of latitude at perihelion \mbox{.\,.\,.\,.\,.\,.\,.\,.\,.\,.\,.\,.\,.} & $169^\circ\;10'\;9''3$ \\
Longitude of the ascending node \mbox{.\,.\,.\,.\,.\,.\,.\,.\,.\,.\,.\,.\,.\,.\,.\,.} & $48\quad44\quad14.7$ \\
Inclination of the orbit \mbox{.\,.\,.\,.\,.\,.\,.\,.\,.\,.\,.\,.\,.\,.\,.\,.\,.\,.\,.\,.\,.} & $106\quad40\quad16.4$ \\
\end{tabular}
\end{center}

\Needspace{17\baselineskip}
These elements represent all the observations that have come to our attention so well that no further improvement can be made, much less any determination of the ellipticity, unless good observations from much later dates at other places become known. For lack of space we can give only a comparison of the elements with the observations made here.

\begin{center}
\begin{tabular}{lr|r|r}
& & \multicolumn{2}{c}{Difference.} \\
& & Right asc. & Declination \\
\hline
January & 30 & $-18''8$ & $+34''9$ \\
February & 3 & $+3.9$ & $-6.0$ \\
& 7 & $-13.1$ & $-36.0$ \\
& 9 & $-0.1$ & $+14.7$ \\
& 10 & $+2.6$ & $+7.3$ \\
& 11 & $-1.2$ & $-3.5$ \\
March & 1 & $-0.3$ & $-0.3$ \\
& 5 & $-5.3$ & $-20.6$ \\
\end{tabular}
\end{center}

\Needspace{11\baselineskip}
This comet stands out from others for its close approach to the Sun and must therefore have reached a very great brightness around the time of perihelion. On the assumption that the comet’s light is borrowed sunlight, Mr. \textsc{von Staudt}’s calculation gave the following brightness values.

\begin{center}
\begin{tabular}{lr|r}
January & 21 & $0.139$ \\
February & 10 & $0.258$ \\
March & 1 & $0.832$ \\
& 5 & $1.253$ \\
\end{tabular}
\end{center}
\Needspace{7\baselineskip}
\begin{center}
\begin{tabular}{lr|r}
March & 18 & $19.588$ \\
& 19 & $32.519$ \\
& 20 & $59.932$ \\
& 21 & $110.610$ \\
\end{tabular}
\end{center}
\Needspace{7\baselineskip}
\begin{center}
\begin{tabular}{lr|r}
March & 22 & $128.139$ \\
& 23 & $81.238$ \\
& 24 & $46.005$ \\
& 25 & $27.980$ \\
\end{tabular}
\end{center}

The figures from 18 March onward apply to the time of culmination at Göttingen. In the first days of March the comet was very readily visible to the naked eye and was almost as bright as a star of the third magnitude. There was therefore every reason to hope that, at its greatest brightness, it could be observed in daylight on the meridian with instruments of great light-gathering power, despite its position near the Sun. At our observatory, however, the overcast sky prevented any attempt from 19 to 23 March. On the 18th and 24th, by contrast, when the sky was clear around the time of the comet’s
\SourceLeaf{444}%
culmination, no trace of the comet could be seen, although its position had been calculated accurately in advance. According to reports that Councillor \textsc{Gauss} has received from several other observatories, unfavorable weather likewise frustrated searches there around the time of the comet’s greatest brightness.
